\chapter{Market making sur ETF}\label{ch:hf:etf-market-making}

Un fonds indiciel coté (ETF) d'obligations se négocie toute la journée en bourse alors que la plupart de ses obligations n'ont pas été échangées depuis des heures. Le teneur de marché qui
en affiche la cotation est, pendant ces heures, le prix. Dans le fonds simulé de ce chapitre, un gel de cinq jours sur le marché obligataire laisse la valeur liquidative officielle
5.6~\% au-dessus de ce que valent les obligations~; un teneur de marché qui affiche autour de ce chiffre perd 70~millions de dollars en cinq jours, soit environ 530 jours de son
revenu habituel. Celui qui valorise le fonds à partir d'un panier déplacé par un proxy liquide gagne 3.2 à 3.6~millions de dollars sur les mêmes jours, et la «~décote~» de 5.3~\% qu'il affiche
à la valeur liquidative n'est presque entièrement due qu'au caractère périmé de celle-ci.

\section{Valorisation d'un ETF à partir de son panier}

Le livre~1, chapitre~14 décrit la structure~: le fonds n'émet et ne rachète ses actions qu'en unités de création, contre un panier, auprès des participants autorisés~;
tous les autres négocient les actions en bourse, où les teneurs de marché les cotent. Le premier problème d'un teneur de marché est un \omterm{def:hf:fair-value:fair}{prix juste} pour les actions
(chapitre~2) lorsque les valeurs publiées par le fonds lui-même sont en retard.

\begin{definition}[Panier de valorisation]\label{def:hf:etf-market-making:pricing}
Un \emph{panier de valorisation}\index{panier de valorisation} est l'ensemble des instruments, des pondérations et des ajustements qu'un teneur de marché utilise pour calculer la fair value d'un ETF
en temps réel~: les titres du fonds là où ils se négocient, un prix modèle là où ils ne se négocient pas (le dernier cours déplacé par un proxy liquide), la devise, les
liquidités et les revenus courus, diminués des frais du fonds.
\end{definition}

La valeur indicative publiée utilise les derniers cours des titres. Pour un fonds dont tous les titres se négocient à chaque seconde, cela suffit. Pour un fonds
d'obligations, ou d'actions dont le marché est fermé, ce n'est pas le cas~: chaque cours ancien n'est qu'une supposition faite à la date de la dernière transaction. Le \omterm{def:hf:etf-market-making:pricing}{panier de valorisation} remplace
chaque cours périmé $P_i$, dont la dernière transaction a eu lieu à $t_i$, par
\[
\hat P_{i,t}=P_i\bigl(1+\beta_i\,(G_t-G_{t_i})\bigr),
\]
où $G$ est un proxy liquide (un contrat à terme sur indice, un indice de crédit, un contrat à terme sur obligations du Trésor) et $\beta_i$ la sensibilité du titre à celui-ci. Pour un fonds dont
le marché sous-jacent est fermé, la même formule s'applique à tous les titres à la fois, avec le mouvement de la devise~: un fonds d'actions japonaises coté à New York
vaut sa clôture de Tokyo multipliée par (1 + le mouvement d'un contrat à terme sur le Nikkei négocié pendant la nuit), puis par (1 + le mouvement du yen).

Petajisto a mesuré à quel point les cours des ETF s'écartent de leurs valeurs liquidatives~: en général dans une bande d'environ 200 points de base, davantage pour les fonds de
titres internationaux ou peu liquides. En contrôlant les cours périmés par ceux de fonds comparables, la bande restait d'environ 100 points de base.

\section{Le marché primaire~: création et rachat}

Le marché primaire est le lieu où l'inventaire d'un teneur de marché peut aller lorsque le marché secondaire refuse de le prendre. Un participant autorisé (ou un
teneur de marché qui passe par lui) livre le panier de création et reçoit une unité de création, ou l'inverse. Deux caractéristiques comptent pour un teneur de marché.

\begin{definition}[Frais de création]\label{def:hf:etf-market-making:fee}
Les \emph{frais de création}\index{frais de création} sont la commission prélevée par le fonds sur chaque création ou rachat d'unité de création~: des frais de transaction fixes par
ordre, et, lorsque le fonds accepte ou verse des espèces au lieu de titres, une commission variable, dans la limite d'un pourcentage annoncé de la valeur de l'unité, pour couvrir les
coûts de négociation du fonds.
\end{definition}

\begin{definition}[Panier sur mesure]\label{def:hf:etf-market-making:custom}
Un \emph{panier sur mesure}\index{panier sur mesure} (custom basket) est un panier de création ou de rachat qui n'est pas une fraction au prorata du portefeuille du fonds, ou qui diffère
du panier publié pour la journée, convenu entre le fonds et le participant autorisé selon les politiques écrites du fonds.
\end{definition}

\begin{dated}{2026-09}{Frais de création, taille des unités et paniers sur mesure}\label{dat:hf:etf-market-making:fees}
Dans sa déclaration d'informations complémentaires datée du 31 juillet 2020 (révisée le 27 novembre 2020), l'iShares Core S\&P 500 ETF fixait à 50\,000 actions par
unité de création, d'une valeur d'environ 14.6~millions de dollars au 30 avril 2020, avec des frais de transaction standard de 1\,250~\$ à la création, et le même montant au rachat, avec des
surcharges maximales de 3.0~\% pour les créations et de 2.0~\% pour les rachats~; les frais standard de l'iShares Europe ETF étaient de 10\,000~\$. La Rule 6c-11 de la SEC,
adoptée en septembre 2019, permet à un ETF qui s'y réfère d'utiliser des \omterm{def:hf:etf-market-making:custom}{paniers sur mesure} s'il adopte des politiques et procédures écrites pour leur construction et
leur acceptation, et exige une transparence quotidienne du portefeuille sur son site. Flow Traders a fait état de 1\,940~milliards d'euros de valeur de produits négociés en bourse (ETP) échangés
en 2025.
\end{dated}

Des frais fixes de 1\,250~\$ sur 14.6~millions de dollars représentent moins d'un point de base~: dans un grand fonds d'actions, les frais ne sont pas la contrainte, c'est le panier. Pour un
fonds obligataire, le panier peut contenir des obligations qu'on ne peut pas acheter aujourd'hui, et le choix du panier par le fonds (livre~9, chapitre~19) détermine quelles obligations un
participant autorisé doit livrer ou recevoir. Le rachat en nature remet au teneur de marché des obligations qu'il doit ensuite vendre, à la
fourchette du marché obligataire~; la décision de création-rachat compare l'écart entre le cours des actions et la valeur du panier, avec les frais et le coût de négociation du panier~:
\[
\text{créer si } (P^{\mathrm{ETF}}-V)\,N>F+c_{\mathrm{buy}}VN,\qquad\text{racheter si } (V-P^{\mathrm{ETF}})\,N>F+c_{\mathrm{sell}}VN,
\]
avec $N$ actions par unité, $F$ les frais et $c$ le coût du panier par dollar. Les décalages de règlement-livraison ajoutent un financement des deux côtés.

\section{Couverture d'un portefeuille d'ETF}

L'inventaire d'un teneur de marché en actions peut être couvert dans le panier (exacte, mais coûteuse), dans un sous-ensemble de celui-ci, dans un contrat à terme sur l'indice, dans un autre fonds sur
le même indice, ou par création ou rachat. L'arithmétique du chapitre~11 s'applique~: le contrat à terme élimine le risque d'indice à faible coût et laisse l'écart
de suivi. La couverture propre aux ETF est la position compensatrice dans un fonds similaire~: deux fonds S\&P 500, ou un fonds et son contrat à terme, se couvrent
l'un l'autre presque parfaitement, et le teneur de marché qui opère sur les deux touche deux fourchettes sur un seul risque.

\section{Lorsque le sous-jacent est fermé ou peu liquide}

Lorsque le sous-jacent ne peut pas être négocié, la cotation du teneur de marché est le marché. Deux catégories de fonds vivent dans cet état~: les fonds dont les titres se négocient dans un
marché fermé (actions internationales cotées à New York) et les fonds dont les titres se négocient rarement (obligations d'entreprises, obligations municipales, prêts).
Le teneur de marché affiche sa cotation à partir de son \omterm{def:hf:etf-market-making:pricing}{panier de valorisation} et couvre sa position avec le proxy~; l'arbitrage création-rachat attend que le sous-jacent rouvre, et
jusque-là, l'écart entre les actions et la valeur liquidative publiée n'est pas un arbitrage mais une prévision.

\begin{omfigure}
\begin{tikzpicture}
\begin{axis}[width=11cm,height=5.4cm,xlabel={\small jour de négociation},ylabel={\small valeur (100 au jour 37)},xmin=37,xmax=50,ymin=89,ymax=102,
  tick label style={font=\footnotesize},grid=major,grid style={gray!25},table/col sep=comma,
  legend style={font=\scriptsize,at={(0.5,-0.3)},anchor=north},legend columns=1]
\fill[gray!12] (axis cs:40,89) rectangle (axis cs:45,102);
\addplot[black,thick] table[x=day,y=true]{figdata/hft/12-etf-market-making/freeze.csv};
\addplot[omDef,thick,dashed] table[x=day,y=nav]{figdata/hft/12-etf-market-making/freeze.csv};
\addplot[omThm,thin] table[x=day,y=mm]{figdata/hft/12-etf-market-making/freeze.csv};
\legend{valeur réelle des obligations,{valeur liquidative des derniers cours},{mid du teneur de marché (panier de valorisation)}}
\end{axis}
\end{tikzpicture}

\omcaption{Le fonds obligataire synthétique lors d'un gel de cinq jours (en grisé~: le facteur recule de 5~\% avec quatre fois sa volatilité, les obligations sont cotées dix fois moins
souvent et leurs fourchettes sont multipliées par six)~: la valeur réelle de ses 100 obligations, la valeur liquidative calculée à partir de leurs derniers cours, et le mid affiché par un
teneur de marché qui valorise à partir d'un panier déplacé par un proxy liquide, toutes les dix minutes. Données~: \texttt{hf\_etf.market}, \texttt{hf\_etf.run}.}\label{fig:hf:etf-market-making:freeze}
\end{omfigure}

\section{Stress~: décotes, valeurs liquidatives périmées et qui les absorbe}

En période de stress, la valeur liquidative publiée d'un fonds peu liquide est périmée par construction~: les obligations qui la feraient bouger ne sont pas négociées. Les actions
se négocient avec une «~décote~» qui tient en partie à l'erreur de la valeur liquidative, et en partie au prix d'une vente à quelqu'un dont le bilan
se remplit. Le livre~9, chapitre~19 sépare les deux pour un fonds obligataire dans une simulation quotidienne et suit les décotes de mars 2020~; ici la question est
celle du teneur de marché. À la pire clôture du gel simulé (la \cref{fig:hf:etf-market-making:freeze}), le teneur de marché qui cote à partir de son \omterm{def:hf:etf-market-making:pricing}{panier de
valorisation} se situait à 532 points de base sous la valeur liquidative. Sur ce total, 557 venaient du caractère périmé de la valeur liquidative (celle-ci se situait à 557 points de base
au-dessus de la valeur réelle des obligations) et 5 venait de la décote du teneur de marché par rapport à cette valeur, c'est-à-dire du prix qu'il demandait pour prendre en charge les actions des vendeurs pris de panique, avec
un bilan capable de les porter. Un teneur de marché qui aurait affiché la valeur liquidative se serait retrouvé à 557 points de base au-dessus de ce qu'il achetait.

\section{Fiches de stratégie}

\begin{strategyfile}{Teneur de marché sur ETF adossé à un panier}\label{strat:hf:etf-market-making:equity}
\sfield{Qui vous paie, et pourquoi} Des investisseurs qui négocient des parts du fonds en bourse et veulent de l'immédiateté pour une taille que les fourchettes propres au panier rendraient coûteuse.
\sfield{Instruments et plateformes} Un fonds actions domestique coté sur sa bourse de cotation et sur d'autres plateformes~; ses composants~; la création et le rachat par
l'intermédiaire d'un participant autorisé.
\sfield{Signal} Le \omterm{def:hf:etf-market-making:pricing}{panier de valorisation} établi à partir des cours en direct des composants (le \omterm{def:hf:fair-value:fair}{prix juste} du chapitre~2), avec les dividendes, la trésorerie et les frais.
\sfield{Taille et exécution} Coter les parts autour de la valeur du panier avec un skew d'inventaire (chapitres~3 à 4)~; couvrir au-delà d'une bande dans les
composants ou dans un contrat à terme (chapitre~11)~; créer ou racheter lorsque l'inventaire atteint une unité et que l'écart paie les frais.
\sfield{Coûts} Fourchettes et frais sur la couverture, \omterm{def:hf:etf-market-making:fee}{frais de création} (moins d'un point de base dans un grand fonds), financement du règlement-livraison.
\sfield{Comment elle meurt} La concurrence~: de nombreux teneurs de marché cotent les mêmes fonds contre les mêmes paniers, et la fourchette revient à celui qui valorise le plus vite et
couvre le moins cher.
\sfield{Horizon, capacité, infrastructure} De quelques secondes à une journée~; capacité tirée du volume du fonds~; paniers en temps réel pour plusieurs centaines de fonds.
\sfield{Backtest honnête} Cours des composants au moment de chaque cotation, et non à la clôture~; heure limite de création et frais du jour.
\sfield{Sources} Livre~1, chapitre~14~; la déclaration d'information additionnelle d'iShares (encadré daté).
\end{strategyfile}

\begin{strategyfile}{Arbitrage création-rachat}\label{strat:hf:etf-market-making:createredeem}
\sfield{Qui vous paie, et pourquoi} Des investisseurs dont les achats ou les ventes éloignent le cours des parts de la valeur du panier plus vite que l'inventaire des teneurs de marché
ne peut l'absorber.
\sfield{Instruments et plateformes} Les parts, les titres du panier, le marché primaire du fonds.
\sfield{Signal} L'écart entre le cours des parts et la valeur exécutable du panier, net des frais et du coût de négociation du panier.
\sfield{Taille et exécution} Acheter le côté bon marché et vendre le côté cher, par tailles d'unités de création~; créer ou racheter à l'heure limite.
\sfield{Coûts} \omterm{def:hf:etf-market-making:fee}{Frais de création} fixes et variables, fourchettes du panier, financement du règlement-livraison, négociation du \omterm{def:hf:etf-market-making:custom}{panier sur mesure}.
\sfield{Comment elle meurt} C'est le mécanisme qui maintient l'écart à un niveau faible~; dans les fonds liquides, l'écart reste dans les coûts presque tout le temps.
\sfield{Horizon, capacité, infrastructure} Une journée~; capacité limitée par l'écart et par le bilan du participant autorisé.
\sfield{Backtest honnête} Cours exécutables du panier, et non la valeur indicative~; heure limite~; droit du fonds de refuser ou de facturer en espèces.
\sfield{Sources} Livre~1, chapitre~14~; livre~9, chapitre~19~; Rule 6c-11 (encadré daté).
\end{strategyfile}

\begin{strategyfile}{Teneur de marché sur ETF obligataire à valeur liquidative périmée}\label{strat:hf:etf-market-making:bond}
\sfield{Qui vous paie, et pourquoi} Des investisseurs qui veulent une exposition obligataire avec l'immédiateté d'une bourse, et qui paient le plus lorsque le marché obligataire est le moins liquide.
\sfield{Instruments et plateformes} Un fonds obligataire coté en bourse~; un proxy liquide (contrat à terme sur le Treasury, indice de crédit)~; les obligations via des dealers et des plateformes.
\sfield{Signal} Un \omterm{def:hf:etf-market-making:pricing}{panier de valorisation}~: la dernière transaction de chaque obligation, déplacée depuis par le proxy et multipliée par sa sensibilité~; la valeur liquidative publiée est une donnée d'entrée,
et non la réponse.
\sfield{Taille et exécution} Coter autour de la valeur du panier, appliquer un skew avec l'inventaire, élargir avec la volatilité du proxy, couvrir l'exposition au taux et au crédit
dans le proxy~; racheter en nature seulement lorsque l'écart des obligations n'est pas trop large.
\sfield{Coûts} La fourchette du proxy, les écarts des obligations au rachat (six fois plus larges dans le gel du tutoriel) et le risque de base entre le proxy et
les obligations.
\sfield{Comment elle meurt} Coter sur une valeur périmée~: dans le tutoriel, un teneur de marché qui cote autour de la valeur liquidative perd 70~millions de dollars lors d'un gel de cinq jours, contre
3.2 à 3.6~millions de dollars gagnés par le teneur de marché adossé au panier.
\sfield{Horizon, capacité, infrastructure} De quelques minutes à quelques jours~; le bilan est la limite~; un service de valorisation obligataire.
\sfield{Backtest honnête} Cotations obligataires telles qu'elles étaient connues à l'époque, base du proxy, et rachat à des prix obligataires exécutables.
\sfield{Sources} Petajisto (2017)~; livre~9, chapitre~19~; le présent chapitre.
\end{strategyfile}

\begin{strategyfile}{Teneur de marché sur ETF international pendant la fermeture du sous-jacent}\label{strat:hf:etf-market-making:international}
\sfield{Qui vous paie, et pourquoi} Des investisseurs qui négocient des marchés étrangers pendant leurs propres heures d'ouverture, lorsque ces marchés sont fermés.
\sfield{Instruments et plateformes} Un fonds d'actions étrangères coté chez soi~; des contrats à terme sur indices qui se négocient la nuit~; la devise.
\sfield{Signal} Le cours de clôture étranger, multiplié par un plus le mouvement du contrat à terme depuis lors, puis par un plus le mouvement de la devise (\texttt{firm.etfmm.closed\_fair}).
\sfield{Taille et exécution} Coter autour de cette valeur, couvrir dans le contrat à terme et dans la devise~; garder l'inventaire jusqu'à l'ouverture du marché étranger, puis
couvrir sur ce marché, ou créer et racheter contre le panier.
\sfield{Coûts} Fourchettes du contrat à terme et de la devise, écart entre le contrat à terme et le panier à l'ouverture, frais.
\sfield{Comment elle meurt} Des nouvelles propres à quelques titres, que le contrat à terme ne reflète pas, et un écart entre le mouvement nocturne du contrat à terme et
l'ouverture du marché.
\sfield{Horizon, capacité, infrastructure} Les heures de non-recouvrement~; capacité tirée de la profondeur des contrats à terme la nuit.
\sfield{Backtest honnête} Le cours du contrat à terme à chaque heure de cotation~; l'ouverture étrangère, et non la clôture, comme sortie.
\sfield{Sources} Petajisto (2017)~: écarts plus grands dans les fonds internationaux.
\end{strategyfile}

\begin{strategyfile}{Teneur de marché sur ETF adossé à des contrats à terme}\label{strat:hf:etf-market-making:futures}
\sfield{Qui vous paie, et pourquoi} Comme dans la stratégie adossée au panier, mais pour des fonds dont les paniers sont chers ou lents à négocier.
\sfield{Instruments et plateformes} Un fonds indiciel et le contrat à terme sur le même indice, ou deux fonds sur un même indice.
\sfield{Signal} Le cours du contrat à terme, ramené par la base juste (dividendes, financement) aux parts du fonds~; le prix relatif des deux fonds.
\sfield{Taille et exécution} Couvrir l'inventaire dans le contrat à terme à chaque fois qu'il bouge d'un contrat (ou au-delà d'une bande)~; coter les deux fonds et compenser le risque.
\sfield{Coûts} La fourchette du contrat à terme et le roll~; l'écart de suivi entre le contrat à terme et le fonds.
\sfield{Comment elle meurt} Les mouvements de la base~: dividendes, financement, roll~; et un fonds qui suit moins bien son indice que ne le fait le contrat à terme.
\sfield{Horizon, capacité, infrastructure} De quelques secondes à quelques jours~; profondeur du contrat à terme~; modèle de base juste.
\sfield{Backtest honnête} Dates de roll et coûts~; dividendes tels qu'annoncés à l'époque.
\sfield{Sources} Chapitre~8 (le contrat à terme précède le fonds)~; chapitre~11 (couverture).
\end{strategyfile}

\section{Tutoriel~: coter un fonds obligataire à travers un gel}\label{tut:hf:etf-market-making:tut}

\begin{tutorial}
\textbf{Objectif.} Coter un fonds obligataire synthétique à partir de sa valeur liquidative ou d'un \omterm{def:hf:etf-market-making:pricing}{panier de valorisation}, couvrir ou non, aplatir ou non, et mesurer ce que chaque
choix rapporte les jours normaux et pendant un gel. \textbf{État final~:} le tableau ci-dessous et les \cref{fig:hf:etf-market-making:freeze,fig:hf:etf-market-making:cumulative}.
\begin{tutsteps}
\item \textbf{Le marché.} \texttt{firm.etfmm.Market}~: 100 obligations de bêtas compris entre 0.6 et 1.4 par rapport à un facteur de volatilité quotidienne de 0.4~\% et de 0.3~\% pour leur propre composante,
chacune publiant une cotation toutes les deux heures en moyenne, à sa valeur plus ou moins une demi-fourchette de 0.1~\%~; un proxy suit le facteur avec une base qui dérive de
0.1~\% par jour. Les jours 40 à 44 sur 60 forment un gel. La valeur liquidative et le \omterm{def:hf:etf-market-making:pricing}{panier de valorisation} sont calculés à partir des cotations publiées.
\omcode{code/firm/etfmm/firm_etfmm.py}{74}{84}{Cotations publiées, valeur liquidative issue de la dernière cotation, et panier de valorisation déplacé par le proxy.}\label{lst:hf:etf-market-making:market}
\item \textbf{Le teneur de marché} cote chaque minute à cinq points de base de chaque côté (quinze pendant le gel) et applique un skew de deux points de base par tranche
de 1~million de dollars, puis dix points de base de plus au-delà de quarante tranches. Les clients négocient une demi-tranche par minute dans chaque sens, davantage lorsque la cotation est plus proche de la valeur réelle~; pendant le gel,
les vendeurs pris de panique en ajoutent deux fois autant et se soucient peu du prix~; un arbitragiste prend toute cotation située à cinq points de base en deçà de la valeur réelle.
\omcode{code/firm/etfmm/firm_etfmm.py}{115}{126}{La cotation, son skew, et la réponse des clients à celle-ci.}\label{lst:hf:etf-market-making:quote}
\item \textbf{Couvrir} dans le proxy (un demi-point de base) ou non~; \textbf{aplatir} en rachetant des unités entières de 100\,000 parts en nature à la clôture
(des frais de 1\,000~\$ et la demi-fourchette des obligations) ou conserver la position.
\item \textbf{Comparer} cinq configurations sur trois marchés simulés (\texttt{hf\_etf.table}).
\end{tutsteps}
\textbf{Pour aller plus loin.} Remplacer les bêtas constants du proxy par des bêtas estimés à partir des cotations passées~; ajouter un second fonds sur les mêmes obligations~; donner au
teneur de marché un bilan plus petit et trouver la taille à partir de laquelle il cesse d'absorber la panique.
\end{tutorial}

\begin{center}
\small
\setlength{\tabcolsep}{4pt}%
\begin{tabular}{@{}>{\raggedright\arraybackslash}p{5.6cm}>{\centering\arraybackslash}p{2.8cm}>{\centering\arraybackslash}p{2.8cm}>{\centering\arraybackslash}p{2.8cm}@{}}
\toprule
cotation à partir de, couverture, aplatissement & jour normal (\$) & son écart type (\$) & gel, cinq jours (millions de \$)\\
\midrule
valeur liquidative, non couverte, conservation & 132\,100 & 24\,300 & $-70.6$\\
valeur liquidative, couverte, conservation & 111\,600 & 16\,000 & $-69.9$\\
\omterm{def:hf:etf-market-making:pricing}{panier de valorisation}, non couvert, conservation & 163\,600 & 13\,500 & 3.17\\
\omterm{def:hf:etf-market-making:pricing}{panier de valorisation}, couvert, conservation & 149\,100 & 9\,300 & 3.55\\
\omterm{def:hf:etf-market-making:pricing}{panier de valorisation}, couvert, rachat & 148\,300 & 9\,400 & 3.26\\
\bottomrule
\end{tabular}
\end{center}

La valeur liquidative se situe en moyenne à 13.7 points de base de la valeur réelle des obligations pendant les minutes normales, et à 214 points de base pendant le gel, lorsque la dernière cotation publiée de l'obligation moyenne
a onze heures de bourse~; le \omterm{def:hf:etf-market-making:pricing}{panier de valorisation} s'écarte de 3.5 et 9.9 points de base. Les jours normaux, le panier rapporte 24~\% de plus, car les clients qui voient le proxy
négocient contre une cotation périmée. Pendant le gel, le teneur de marché qui cote sur la valeur liquidative achète aux vendeurs pris de panique à des prix que les obligations ne soutiennent plus, et
perd 70~millions de dollars, dont 95~\% sur les transactions elles-mêmes plutôt que chez les arbitragistes. La couverture coûte environ 14\,500~\$ un jour normal et réduit
d'un tiers l'écart type quotidien. Racheter en nature pendant le gel coûte 262\,000~\$ en écarts obligataires, pour un inventaire que le portefeuille couvert aurait pu tenir.

\begin{omfigure}
\begin{tikzpicture}
\begin{axis}[width=11cm,height=5.2cm,xlabel={\small jour de négociation},ylabel={\small P\&L cumulé (millions de \$)},xmin=0,xmax=60,ymin=-95,ymax=20,
  tick label style={font=\footnotesize},grid=major,grid style={gray!25},table/col sep=comma,
  legend style={font=\scriptsize,at={(0.03,0.05)},anchor=south west}]
\fill[gray!12] (axis cs:40,-95) rectangle (axis cs:45,20);
\addplot[omDef,thick,dashed] table[x=day,y=nav]{figdata/hft/12-etf-market-making/cumulative.csv};
\addplot[omThm,thick] table[x=day,y=basket]{figdata/hft/12-etf-market-making/cumulative.csv};
\legend{{cotation sur la valeur liquidative, couverte},{cotation sur le panier de valorisation, couverte}}
\end{axis}
\end{tikzpicture}

\omcaption{P\&L cumulé, valorisé à la valeur réelle des obligations, des deux teneurs de marché couverts sur le premier marché simulé~; le gel est grisé.
Données~: \texttt{hf\_etf.run}.}\label{fig:hf:etf-market-making:cumulative}
\end{omfigure}

\section{Construction~: le teneur de marché de l'ETF}\label{bld:hf:etf-market-making:etfmm}

\begin{build}
\bfield{Rôle} Valoriser un ETF à partir d'un \omterm{def:hf:etf-market-making:pricing}{panier de valorisation} dont certains composants sont périmés et ont un proxy, le coter et le couvrir, et décider des créations et des rachats.
\bfield{Interface} \texttt{closed\_fair(last, beta, proxy\_move, fx\_move)}, \texttt{pricing\_basket(last, t\_last, beta, G, t)},
\texttt{creation\_decision(price, value, unit, fee, cost\_bp)}, \texttt{Market(seed, days, freeze)} avec \texttt{true}, \texttt{nav},
\texttt{basket}, \texttt{age}~; \texttt{run\_mm(market, estimator, hedge, flatten)} renvoyant le P\&L quotidien par composante et le mid coté. Construit
en parallèle de \texttt{firm.etfmonitor} (livre~1), qui affiche la prime à partir des cotations en direct.
\bfield{Règles} Le P\&L est valorisé à la valeur réelle, jamais à la valeur liquidative~; les créations et les rachats portent sur des unités entières~; la couverture s'effectue dans le
proxy, à son propre coût.
\bfield{Tests d'acceptation} \texttt{code/firm/etfmm/tests/}~: valeurs à marché fermé et valeurs du panier calculées à la main~; cas de création, de rachat et d'absence d'action~; la valeur liquidative
est égale à la moyenne des dernières cotations, et chaque dernière cotation est antérieure ou égale à son instant~; l'erreur du panier pendant un gel est inférieure à un cinquième de celle de la
valeur liquidative~; les composantes du P\&L s'additionnent au total~; le teneur de marché qui cote sur le panier surpasse celui qui cote sur la valeur liquidative pendant le gel~; la couverture réduit la
variabilité du P\&L d'inventaire.
\bfield{Pour aller plus loin} Bêtas estimés~; \omterm{def:hf:etf-market-making:custom}{paniers sur mesure} choisis par le fonds~; un second fonds sur les mêmes obligations~; une limite de bilan qui devient contraignante.
\end{build}

\begin{omsources}
\item A.~Petajisto, Inefficiencies in the pricing of exchange-traded funds, \emph{Financial Analysts Journal} 73(1), 2017, 24--54.
\item iShares Trust, Statement of Additional Information du 31 juillet 2020 (version révisée le 27 novembre 2020), SEC Form 497.
\item US Securities and Exchange Commission, communiqué de presse 2019-190 sur la Rule 6c-11, 26 septembre 2019.
\item Flow Traders, résultats du T4 et de l'exercice 2025, 12 février 2026.
\end{omsources}

\section{Exercices}

\begin{exercise}[$\star$]\label{exo:hf:etf-market-making:1}
Exprimer des frais de 1\,250~\$ sur une unité de création de 50\,000 actions à 291.58~\$ en points de base.
\end{exercise}

\begin{exercise}[$\star$]\label{exo:hf:etf-market-making:2}
Le \omterm{def:hf:etf-market-making:pricing}{panier de valorisation} d'un fonds d'actions japonaises valait 100 à la clôture. Depuis, un contrat à terme sur le Nikkei a progressé de 1.5~\% et le yen a perdu 0.5~\% face au
dollar. Avec un bêta égal à un, quelle est la juste valeur du fonds~?
\end{exercise}

\begin{exercise}[$\star$]\label{exo:hf:etf-market-making:3}
Les actions s'échangent à 100.05, le panier vaut 100.00, une unité compte 50\,000 actions, les frais s'élèvent à 1\,250~\$ et l'achat du panier coûte 2 points de base. Créer,
racheter ou ne rien faire, et avec quel avantage~?
\end{exercise}

\begin{exercise}[$\star\star$]\label{exo:hf:etf-market-making:4}
À la clôture la plus basse du gel, le mid du teneur de marché se situait 532 points de base sous la valeur liquidative. Décomposer cet écart entre le caractère périmé et la
décote propre au teneur de marché, et dire lequel un arbitragiste pourrait négocier.
\end{exercise}

\begin{exercise}[$\star\star$]\label{exo:hf:etf-market-making:5}
Pourquoi le teneur de marché qui cote sur la valeur liquidative perd-il, même les jours normaux, par rapport à celui qui cote sur le panier~?
\end{exercise}

\begin{exercise}[$\star\star$]\label{exo:hf:etf-market-making:6}
Pourquoi le rachat en nature pendant le gel a-t-il diminué le P\&L du teneur de marché couvert~?
\end{exercise}

\begin{exercise}[$\star\star\star$]\label{exo:hf:etf-market-making:7}
\emph{Programmation.} Faire publier les obligations cent fois moins souvent pendant le gel au lieu de dix fois (\texttt{hf\_etf.sparser}). Que deviennent le caractère périmé maximal de la
valeur liquidative, l'erreur du \omterm{def:hf:etf-market-making:pricing}{panier de valorisation} et le P\&L des deux teneurs de marché couverts pendant le gel du premier marché~?
\end{exercise}

\begin{exercise}[$\star\star\star$]\label{exo:hf:etf-market-making:8}
\emph{Trouver la faille.} «~Le fonds a échangé 5~\% sous sa valeur liquidative, donc nous avons acheté les actions et les avons rachetées~: un gain sans risque de 5~\%.~»
\end{exercise}

\section{Problème~: cotation d'un marché fermé}

\begin{problem}[{Problème du week-end --- cotation d'un marché fermé}]\label{pb:hf:etf-market-making:1}
Un teneur de marché cote un fonds obligataire pendant un gel de cinq jours du marché obligataire.

\medskip\noindent\textbf{Partie I --- Prix.}
\begin{enumerate}
\item Définir un \omterm{def:hf:etf-market-making:pricing}{panier de valorisation} et écrire l'ajustement de prix périmé.
\item Comment le même ajustement valorise-t-il un fonds international pendant la fermeture de son marché~?
\item Qu'a établi Petajisto sur les prix des ETF et les valeurs liquidatives~?
\item Dans quelle mesure la valeur liquidative et le \omterm{def:hf:etf-market-making:pricing}{panier de valorisation} étaient-ils périmés pendant les minutes normales et pendant le gel~?
\end{enumerate}

\medskip\noindent\textbf{Partie II --- Le marché primaire.}
\begin{enumerate}[resume]
\item Définir des \omterm{def:hf:etf-market-making:fee}{frais de création} et donner l'exemple de l'encadré daté.
\item Définir un \omterm{def:hf:etf-market-making:custom}{panier sur mesure} et la règle qui le permet.
\item Écrire la décision de création-rachat.
\item Pourquoi le rachat en nature d'un fonds obligataire n'est-il pas gratuit pendant un gel~?
\end{enumerate}

\medskip\noindent\textbf{Partie III --- La simulation.}
\begin{enumerate}[resume]
\item Donner le P\&L d'une journée normale et son écart type pour les teneurs de marché qui cotent sur la valeur liquidative et sur le panier, couverts.
\item Que chacun a-t-il perdu ou gagné pendant le gel, et d'où provient la perte du teneur de marché qui cote sur la valeur liquidative~?
\item Combien a coûté la couverture et qu'a-t-elle permis d'acheter~?
\item Combien a coûté le rachat pendant le gel~?
\end{enumerate}

\medskip\noindent\textbf{Partie IV --- Le verdict.}
\begin{enumerate}[resume]
\item Énoncer le \emph{résultat nommé}~: le P\&L du teneur de marché et la prime qu'il a cotée pendant le gel, et la part de la décote due au caractère périmé de la
valeur liquidative.
\item Qui a absorbé les actions des vendeurs paniqués, et à quel prix~?
\item Que changerait un bilan plus petit~?
\item Quand une décote sur la valeur liquidative constitue-t-elle un arbitrage~?
\item Que doivent porter les valeurs publiées d'un fonds pour aider les investisseurs pendant un gel~?
\item Comment estimer les bêtas des obligations par rapport au proxy~?
\item Laquelle des cinq fiches de stratégie est la plus exposée à un gel, et pourquoi~?
\item En une phrase~: quel est le prix d'un teneur de marché d'ETF sur un marché fermé~?
\end{enumerate}
\end{problem}

\section{Questions d'entretien}

\begin{interviewq}[$\star$ \iqroles{trader}]\label{iq:hf:etf-market-making:1}
Un fonds S\&P 500 vous est proposé à un cent sous la valeur de son panier. Que faites-vous, et qu'est-ce qui empêche tout le monde d'en faire autant~?
\end{interviewq}

\begin{interviewq}[$\star\star$ \iqroles{chercheur}]\label{iq:hf:etf-market-making:2}
Comment valoriseriez-vous en cours de séance un fonds d'obligations à haut rendement quand la moitié de ses obligations n'ont pas été négociées aujourd'hui~?
\end{interviewq}

\begin{interviewq}[$\star\star$ \iqroles{trader}]\label{iq:hf:etf-market-making:3}
Vous détenez une position acheteuse de 50~millions de dollars d'un fonds obligataire à 15:55 dans une phase de baisse. Racheter, couvrir ou conserver~? Que devez-vous savoir~?
\end{interviewq}

\begin{interviewq}[$\star\star$ \iqroles{risque}]\label{iq:hf:etf-market-making:4}
Votre P\&L valorise le fonds à sa valeur liquidative. Pourquoi est-ce dangereux, et à quoi devrait-il être valorisé~?
\end{interviewq}

\begin{interviewq}[$\star\star$ \iqroles{développeur}]\label{iq:hf:etf-market-making:5}
Concevez le service qui calcule les \omterm{def:hf:etf-market-making:pricing}{paniers de valorisation} de 2\,000 fonds en temps réel.
\end{interviewq}

\begin{interviewq}[$\star\star\star$ \iqroles{chercheur}]\label{iq:hf:etf-market-making:6}
Montrez que l'ajustement de prix périmé est la meilleure prévision linéaire de la valeur d'une détention à partir de son dernier prix et du mouvement du proxy, et dites quand elle est
biaisée.
\end{interviewq}
